% Generated draft: requires mathematical and source review.
\documentclass[11pt]{article}
\usepackage[margin=25mm]{geometry}
\usepackage{fontspec}
\setmainfont{Latin Modern Roman}
\usepackage{amsmath,amssymb,mathtools,unicode-math}
\setmathfont{Latin Modern Math}
\usepackage{enumitem,needspace,xurl,hyperref}
\setlength{\parindent}{0pt}
\setlength{\parskip}{5pt}
\setlength{\emergencystretch}{3em}
\newcommand{\N}{\mathbb{N}}
\newcommand{\Z}{\mathbb{Z}}
\newcommand{\Q}{\mathbb{Q}}
\newcommand{\R}{\mathbb{R}}
\newcommand{\C}{\mathbb{C}}
\newcommand{\F}{\mathbb{F}}
\newcommand{\A}{\mathbb{A}}
\newcommand{\PP}{\mathbb{P}}
\newcommand{\E}{\mathbb{E}}
\newcommand{\eps}{\varepsilon}
\newcommand{\dd}{\,\mathrm{d}}
\begin{document}
\section*{UnsolvedMath ID: 20002698}
\textbf{Source problem code:} catalogue entry
\subsection{Definitions and mathematical statement}
The canonical fragment is not an independent Problem 3.2. The aid identifies it as the latter half of a remarks bullet after Problem 1.1, with an internal reference to Nathan Reading's later four-part sortable-elements Problem 3.2. Its actual content says that Tom Brady and Colum Watt gave a new definition of $\Delta_W$ using noncrossing partitions, reference [14], suggesting a structural connection between $NC_W$ and $\Delta_W$. The referenced four questions and bibliography entry [14] are not included. No separate quantified problem is stated by this fragment.

The aid develops a finite Coxeter-system compatibility statement. Let $(W,S)$ be finite, fix a Coxeter element $c$, and let $\Delta_c(W)$ be its labeled $c$-cluster complex on almost-positive roots. With $T$ the set of reflections, define
\[
x\leq_Ty\ \Longleftrightarrow\ \ell_T(x)+\ell_T(x^{-1}y)=\ell_T(y),\qquad
NC_c(W)=[1,c]_{\leq_T}.
\]
Use Reading's bijections $\operatorname{cl}_c$ from $c$-sortable elements to cluster facets and $\operatorname{nc}_c$ from those elements to $NC_c(W)$, and put
\[
\Phi_c=\operatorname{nc}_c\circ\operatorname{cl}_c^{-1}.
\]
For $I\subseteq S$, let $J=S\setminus I$, $W_J=\langle J\rangle$, $c_J$ be obtained by deleting letters outside $J$ from a Coxeter word for $c$, and $F_I=\{-\alpha_s:s\in I\}$. The supplied theorem is
\[
\operatorname{link}_{\Delta_c(W)}(F_I)\cong\Delta_{c_J}(W_J),
\qquad
F_I\subseteq C\ \Longleftrightarrow\ \Phi_c(C)\in W_J
\]
for every cluster facet $C$. Moreover $\Phi_c$ restricts to a bijection
\[
\{C:F_I\subseteq C\}\longrightarrow
NC_c(W)\cap W_J=NC_{c_J}(W_J),
\]
compatible with the parabolic Reading construction after removing the fixed face from a facet. It includes $I=\varnothing$ and $I=S$ with their usual empty-parabolic conventions.

The aid calls this restriction rank preserving, but facets do not carry the varying face-poset ranks of noncrossing partitions. The appropriate statistic must be specified: with the sortable-element interpretation, use the number of right descents of $\operatorname{cl}_c^{-1}(C)$, corresponding to $\ell_T(\Phi_c(C))$. Parabolic restriction preserves that statistic. No isomorphism of the entire cluster face poset with the noncrossing-partition poset, or an undefined varying rank on facets, is asserted.

The summary reports that the original suggested finite connection was subsequently realized through Reading's sortable bijections, Reading--Speyer uniform arguments, and Brady--Watt facet/noncrossing and Cambrian-fan results. Those are supplied context claims without full references. The displayed compatibility theorem is a developed synthesis, not an open question contained in the extracted remarks fragment.
\subsection{Short English statement}
The source is a context fragment pointing to a connection between noncrossing partitions and cluster complexes, not an independent problem. The aid states a finite Coxeter-system bridge: fixing negative-simple roots in a cluster corresponds to restricting its noncrossing image to the complementary parabolic subgroup. Its rank claim requires the sortable-descent statistic rather than ordinary facet dimension.
\subsection{Sources}
\begin{itemize}
\item Local UnsolvedMath registry, permanent ID $20002698$, \texttt{catalogue entry}; \texttt{aim-probability-notes.json}, \texttt{notes[139]}, extracted item ``3.2.''
\item American Institute of Mathematics, \emph{Braid groups, clusters and free probability}, \url{https://aimath.org/WWN/braidgroups/braidgroups.pdf}. The aid identifies the passage as remarks following Problem 1.1 and its Problem 3.2 reference as internal to Nathan Reading's later sortable-elements problem.
\item Original attribution: Tom Brady and Colum Watt, reference [14]. The full entry is missing. Supplied context additionally names Reading, Reading--Speyer, and Brady--Watt for sortable, uniform, facet/noncrossing and Cambrian-fan connections; no authoritative full bibliographic details are supplied.
\item Research attempt 1: context only, status ``not a problem'' at run, rigor label ``proved,'' novelty confidence low. Registry labels open and open triage, statement status reconstructed/unverified, do not make this fragment an independent conjecture.
\end{itemize}
\end{document}
